\makeatletter
\@ifundefined{isChecklistMainFile}{
  \newif\ifreproStandalone
  \reproStandalonetrue
}{
  \newif\ifreproStandalone
  \reproStandalonefalse
}
\makeatother

\ifreproStandalone
\documentclass[letterpaper]{article}
\usepackage[submission]{aaai2027}
\frenchspacing
\begin{document}
\fi

\setlength{\leftmargini}{20pt}
\makeatletter
\def\@listi{\leftmargin\leftmargini \topsep .5em \parsep .5em \itemsep .5em}
\def\@listii{\leftmargin\leftmarginii \labelwidth\leftmarginii \advance\labelwidth-\labelsep \topsep .4em \parsep .4em \itemsep .4em}
\def\@listiii{\leftmargin\leftmarginiii \labelwidth\leftmarginiii \advance\labelwidth-\labelsep \topsep .4em \parsep .4em \itemsep .4em}
\makeatother

\setcounter{secnumdepth}{0}
\renewcommand\thesubsection{\arabic{subsection}}
\renewcommand\labelenumi{\thesubsection.\arabic{enumi}}

\newcounter{checksubsection}
\newcounter{checkitem}[checksubsection]

\newcommand{\checksubsection}[1]{%
  \refstepcounter{checksubsection}%
  \paragraph{\arabic{checksubsection}. #1}%
  \setcounter{checkitem}{0}%
}

\newcommand{\checkitem}{%
  \refstepcounter{checkitem}%
  \item[\arabic{checksubsection}.\arabic{checkitem}.]%
}
\newcommand{\question}[2]{\normalcolor\checkitem #1 #2 \color{blue}}
\newcommand{\ifyespoints}[1]{\makebox[0pt][l]{\hspace{-15pt}\normalcolor #1}}

\section*{Reproducibility Checklist}

This checklist follows the official AAAI-27 Author Kit and public AAAI-27
main-track instructions checked on 2026-06-13. Answers refer to the main paper
plus the Docker-based experiment artifact bundle planned for anonymous
OpenReview supplementary upload if allowed by the portal, and for
post-acceptance GitHub/Zenodo release.

\checksubsection{General Paper Structure}
\begin{itemize}

\question{Includes a conceptual outline and/or pseudocode description of AI methods introduced}{(yes/partial/no/NA)}
yes. Sections 3--4 and Figure 1 define the row contract, geometry join,
calibrated geometry-consistency score, and re-ranking/evaluation protocol.

\question{Clearly delineates statements that are opinions, hypothesis, and speculation from objective facts and results}{(yes/no)}
yes. The paper states the scoped reliability claim, validation scope,
Open3DSG recovery-policy caveat, blocked extensions, and limitations
separately from empirical results.

\question{Provides well-marked pedagogical references for less-familiar readers to gain background necessary to replicate the paper}{(yes/no)}
yes. Related Work cites 3D scene graph prediction, open-vocabulary 3D scene
graph methods, geometry-aware relation modeling, and reliability/calibration
work.

\end{itemize}
\checksubsection{Theoretical Contributions}
\begin{itemize}

\question{Does this paper make theoretical contributions?}{(yes/no)}
no. The contribution is methodological and empirical.

\ifyespoints{\vspace{1.2em}If yes, please address the following points:}
\begin{itemize}

\question{All assumptions and restrictions are stated clearly and formally}{(yes/partial/no)}
NA.

\question{All novel claims are stated formally (e.g., in theorem statements)}{(yes/partial/no)}
NA.

\question{Proofs of all novel claims are included}{(yes/partial/no)}
NA.

\question{Proof sketches or intuitions are given for complex and/or novel results}{(yes/partial/no)}
NA.

\question{Appropriate citations to theoretical tools used are given}{(yes/partial/no)}
NA.

\question{All theoretical claims are demonstrated empirically to hold}{(yes/partial/no/NA)}
NA.

\question{All experimental code used to eliminate or disprove claims is included}{(yes/no/NA)}
NA.

\end{itemize}
\end{itemize}

\checksubsection{Dataset Usage}
\begin{itemize}

\question{Does this paper rely on one or more datasets?}{(yes/no)}
yes.

\ifyespoints{If yes, please address the following points:}
\begin{itemize}

\question{A motivation is given for why the experiments are conducted on the selected datasets}{(yes/partial/no/NA)}
yes. Open3DSG is the main open-vocabulary relation-source case study and
VL-SAT is the controlled reproduced anchor under the same row contract.

\question{All novel datasets introduced in this paper are included in a data appendix}{(yes/partial/no/NA)}
NA. The paper introduces no new dataset.

\question{All novel datasets introduced in this paper will be made publicly available upon publication of the paper with a license that allows free usage for research purposes}{(yes/partial/no/NA)}
NA.

\question{All datasets drawn from the existing literature (potentially including authors' own previously published work) are accompanied by appropriate citations}{(yes/no/NA)}
yes. The paper cites the 3DSSG, 3RScan, VL-SAT, and Open3DSG sources.

\question{All datasets drawn from the existing literature (potentially including authors' own previously published work) are publicly available}{(yes/partial/no/NA)}
partial. The sources are public research routes, but raw scans, meshes,
feature caches, checkpoints, and model caches remain subject to the original
access terms and are not redistributed in the PDF.

\question{All datasets that are not publicly available are described in detail, with explanation why publicly available alternatives are not scientifically satisficing}{(yes/partial/no/NA)}
NA. No private dataset is introduced.

\end{itemize}
\end{itemize}

\checksubsection{Computational Experiments}
\begin{itemize}

\question{Does this paper include computational experiments?}{(yes/no)}
yes.

\ifyespoints{If yes, please address the following points:}
\begin{itemize}

\question{This paper states the number and range of values tried per (hyper-) parameter during development of the paper, along with the criterion used for selecting the final parameter setting}{(yes/partial/no/NA)}
partial. The paper reports the operating points, family policies, exact-label
denominator, source scope, and Open3DSG checkpoint selection by train-dev
validation loss. It does not claim exhaustive hyperparameter sweeps.

\question{Any code required for pre-processing data is included in the appendix}{(yes/partial/no)}
partial. Docker runbooks and scripts cover preprocessing, feature generation,
adapter export, geometry join, metrics, tables, and failure analysis; large
runtime artifacts are rebuilt, mounted, or released separately.

\question{All source code required for conducting and analyzing the experiments is included in a code appendix}{(yes/partial/no)}
partial. The tracked experiment root contains Docker entrypoints, scripts,
manifests, and reports. Large datasets, model caches, feature caches, and
selected checkpoints are external artifacts.

\question{All source code required for conducting and analyzing the experiments will be made publicly available upon publication of the paper with a license that allows free usage for research purposes}{(yes/partial/no)}
partial. The release plan is fixed as GitHub source code plus a Zenodo DOI for
the fixed full-validation result bundle after anonymization is no longer
required, but the final release URL/DOI and license files are not yet fixed.

\question{All source code implementing new methods have comments detailing the implementation, with references to the paper where each step comes from}{(yes/partial/no)}
partial. The scripts and runbooks document the row contract, source adapters,
geometry scoring, calibration, and metric generation; final code-release
cleanup remains.

\question{If an algorithm depends on randomness, then the method used for setting seeds is described in a way sufficient to allow replication of results}{(yes/partial/no/NA)}
partial. Evaluation is deterministic after fixed predictions, checkpoints,
and manifests; reproduced source training randomness is recorded as
provenance rather than evaluated through a multi-seed study.

\question{This paper specifies the computing infrastructure used for running experiments (hardware and software), including GPU/CPU models; amount of memory; operating system; names and versions of relevant software libraries and frameworks}{(yes/partial/no)}
yes. The Docker runbooks and artifact bundle record the exact compose routes,
mounted paths, and image ids. The verification host uses Ubuntu 24.04.4 LTS,
62 GiB RAM, an NVIDIA GeForce RTX 5090 with 32607 MiB GPU memory, driver
580.159.03, Docker 29.5.2, Docker Compose v5.1.4, and Python 3.12.3. The
tracked hardware/software manifest records image ids for
\texttt{h001-geom-reliability:latest}, \texttt{h001-open3dsg-repro:cu128},
and \texttt{h001-aaai-tex:20260611}.

\question{This paper formally describes evaluation metrics used and explains the motivation for choosing these metrics}{(yes/partial/no)}
yes. The paper defines exact-label \rAt{50}/\rAt{100},
\vAt{50}/\vAt{100}, GT/counterfactual verifier metrics, controls, bootstrap
CI, and failure-analysis criteria.

\question{This paper states the number of algorithm runs used to compute each reported result}{(yes/no)}
yes. Reported source-result tables use one fixed reproduced source run per
source/condition under locked manifests.

\question{Analysis of experiments goes beyond single-dimensional summaries of performance (e.g., average; median) to include measures of variation, confidence, or other distributional information}{(yes/no)}
yes. The paper reports recall/violation tradeoffs, source/family operating
points, counterfactual verifier checks, controls, failure rows, qualitative
case inspection, and Docker subgraph bootstrap confidence intervals.

\question{The significance of any improvement or decrease in performance is judged using appropriate statistical tests (e.g., Wilcoxon signed-rank)}{(yes/partial/no)}
partial. The paper uses paired subgraph bootstrap confidence intervals for
source-result deltas. It does not report repeated-run variance or Wilcoxon
tests.

\question{This paper lists all final (hyper-)parameters used for each model/algorithm in the paper's experiments}{(yes/partial/no/NA)}
partial. The paper and artifact reports list operating points, checkpoint
choice, denominator policy, and source caveats, including the Open3DSG
548/548 recovery branch with \texttt{OPEN3DSG\_MIN\_VISIBLE\_OBJECTS=2} and
relaxed two-scan view regeneration. A final release table should collect every
threshold and configuration in one manifest.

\end{itemize}
\end{itemize}

\ifreproStandalone
\end{document}
\fi
